% LSPV-5 UNIVERSALDAG1 — 2026-10-02
\section*{UNIVERSALDAG1: Exact finite universal solver-DAG baseline for small 3-SAT}

\textbf{STATUS: EXACT FINITE BASELINE / CANDIDATE SCALING PROGRAM.}
This experiment does not prove a SAT lower bound or separate P from NP.

\subsection*{Literature audit}
The proposed experiment overlaps substantially with established branching-program and knowledge-compilation theory, but the exact finite all-formulas coverage experiment remains useful as a controlled laboratory.

For a fixed variable order, reduced OBDDs are canonical and minimal up to isomorphism, and can be reduced efficiently. Exact minimization of free BDDs has also been studied for small Boolean functions. Huang and Darwiche showed that DPLL augmented with caching can construct OBDDs directly from CNF, with bounds exponential in cutwidth/pathwidth of the variable ordering with respect to a hypergraph abstraction of the CNF. Decision-DNNF/FBDD work supplies further lower-bound and representation controls. Therefore any claim that a large explicit DPLL tree or DAG is intrinsically necessary must be compared against OBDD, FBDD, decision-DNNF and other compiled representations.

Relevant references:
Sieling and Wegener, Reduction of OBDDs in linear time, Information Processing Letters 48 (1993), 139--144.
Günther and Drechsler, Minimization of free BDDs, Integration 32 (2002), 41--59.
Huang and Darwiche, Using DPLL for Efficient OBDD Construction, SAT 2004.
Beame, Li, Roy and Suciu, Lower Bounds for Exact Model Counting and Applications in Probabilistic Databases, 2013/2014.
Samar, Vinyals and Ganesh, An Exponential Separation Between Deterministic CDCL and DPLL Solvers, SAT 2026.

\subsection*{Model correction}
The initially proposed n=3, m=8 case with simple full 3-clauses is degenerate. There are only eight possible signed clauses and each excludes exactly one of the eight assignments. Selecting all eight is the unique 8-clause set and is UNSAT.

The first useful exhaustive boundary is n=4, m=8. There are
\[
8\binom43=32
\]
possible simple full 3-clauses and
\[
\binom{32}{8}=10,518,300
\]
distinct 8-clause formulas.

\subsection*{Exact cube interpretation}
For n=4, each 3-clause omits one variable and is falsified by exactly two assignments differing only in that omitted variable. Hence the 32 possible clauses correspond exactly to the 32 edges of the Boolean 4-cube Q_4.

An 8-clause formula corresponds to an 8-edge subset E of Q_4. Its satisfying assignments are precisely the cube vertices not covered by E. Therefore
\[
F_E\text{ is UNSAT}\iff E\text{ covers all 16 vertices of }Q_4.
\]
This gives an exact semantic quotient and an immediate hostile compact representation.

\subsection*{Exhaustive results}
The complete set of 10,518,300 formulas was enumerated. It produced 29,417 distinct semantic falsified-assignment masks and exactly 272 UNSAT formulas.

The number of distinct semantic masks by number of covered cube vertices was:
\[
\begin{array}{c|rrrrrrrrrr}
\text{covered}&7&8&9&10&11&12&13&14&15&16\\
\hline
\text{masks}&688&5852&9008&7176&4192&1804&560&120&16&1.
\end{array}
\]

Thus more than ten million syntactically distinct formulas collapse to fewer than thirty thousand semantic states in this finite universe.

\subsection*{Meaning for the universal-DAG idea}
This is already a strong warning against counting formulas or branches. The universal solver may quotient many syntactic formulas into the same semantic state. The experimental object should therefore be a ladder of increasingly powerful quotients:
\[
\text{syntactic formulas}
\to
\text{hypergraph-isomorphism classes}
\to
\text{DPLL residual states}
\to
\text{reduced OBDD/FBDD states}
\to
\text{semantic assignment masks}.
\]
At each stage record the number of states, edges, maximum width, binary log of state count, and cost of computing the quotient/selector.

\subsection*{Next exact experiment}
Construct a universal residual-state DAG for n=4,m=8 under explicitly declared DPLL operations:
variable assignment, unit propagation, tautology deletion, satisfied-clause deletion, falsified-literal deletion, duplicate-clause deletion, and optional subsumption in separate model variants.

Canonicalize residual formulas first syntactically and then under legal variable renaming/sign symmetries. Compare the resulting DAG against:
(1) the semantic 16-bit mask quotient;
(2) minimum reduced OBDDs over all 4! variable orders for each semantic function;
(3) an exact FBDD baseline where feasible;
(4) the direct cube-edge-cover representation.

The key quantity is not the number of individual short paths but the size and selector cost of one common representation covering every formula.

\subsection*{Strict guard}
Even an exponential minimum DAG in one declared representation would be only a restricted-model lower bound. An unrestricted polynomial-time SAT algorithm need not instantiate that DAG. Any eventual P-versus-NP implication still requires the arbitrary-algorithm acquisition/necessity bridge and the nonrelativizing/nonalgebrizing hinge.

P versus NP remains OPEN.

% END_OF_NOTE::UNIVERSALDAG1
